\documentclass[aspectratio=169,10pt]{beamer}

\usepackage[portuguese]{babel}
\usepackage[utf8]{inputenc}
\usepackage[T1]{fontenc}
\usepackage{amsmath,amssymb}
\usepackage{booktabs}
\usepackage{graphicx}
\usepackage{tikz}
\usepackage{hyperref}
\usepackage[numbers]{natbib}
\usepackage{braket}

\usetikzlibrary{arrows.meta,positioning,shapes.geometric,calc}

\definecolor{DeepBlue}{HTML}{102A43}
\definecolor{QuantumPurple}{HTML}{6D28D9}
\definecolor{NeuralOrange}{HTML}{F97316}
\definecolor{SoftGray}{HTML}{F3F4F6}
\definecolor{AccentGreen}{HTML}{059669}

\usetheme{Madrid}
\usecolortheme{default}
\setbeamertemplate{navigation symbols}{}
\setbeamertemplate{itemize item}{\color{QuantumPurple}$\blacktriangleright$}
\setbeamertemplate{itemize subitem}{\color{NeuralOrange}$\bullet$}

\setbeamercolor{title}{fg=white,bg=DeepBlue}
\setbeamercolor{frametitle}{fg=white,bg=DeepBlue}
\setbeamercolor{block title}{fg=white,bg=QuantumPurple}
\setbeamercolor{block body}{fg=black,bg=SoftGray}
\setbeamercolor{alerted text}{fg=NeuralOrange}

\title[Classificacao Quantica de Imagens]{Classificacao Quantica de Imagens com Extratores Profundos de Caracteristicas}
\subtitle{Proposta de investigacao: codificacao de estados, numero de qubits e classificadores quanticos variacionais}
\author[Lorran de Araujo]{Lorran de Araujo Duraes Soares}
\institute[LNCC]{Topicos Avancados de Redes Neurais Profundas e Aplicacoes\\Laboratorio Nacional de Computacao Cientifica}
\date{\today}

\newcommand{\highlight}[1]{\textcolor{QuantumPurple}{\textbf{#1}}}

\begin{document}

\begin{frame}
  \titlepage
\end{frame}

\begin{frame}{Contextualização}
  \begin{columns}[T,onlytextwidth]
    \begin{column}{0.48\textwidth}
      \begin{block}{Modelos clássicos atuais}
        \begin{itemize}
          \item Redes neurais convolucionais tem apresentado excelente desempenho em classificação de imagens.
          \item Entretanto, modelos totalmente clássicos podem exigir muitos parâmetros e grande capacidade computacional.
        \end{itemize}
      \end{block}
    \end{column}

    \begin{column}{0.48\textwidth}
      \begin{block}{Motivação quântica}
        \begin{itemize}
          \item Circuitos quânticos exploram propriedades como superposição, interferência e emaranhamento; 
          \item Ofere uma forma alternativa de representar e processar dados de alta dimensão.
        \end{itemize}
      \end{block}
    \end{column}
  \end{columns}

  \vspace{0.35cm}

  \begin{alertblock}{Questão de interesse}
    No cenário atual, uma abordagem prática é utilizar arquiteturas híbridas:
    
    \begin{itemize}
      \item uma CNN clássica extrai as características;
      \item um circuito quântico variacional realiza a classificação.
    \end{itemize}

  \end{alertblock}
\end{frame}

\begin{frame}{Ideia Central}
  \begin{block}{Proposta}
    Desenvolver uma arquitetura híbrida composta por um extrator de características clássico e por circuitos quânticos variacionais utilizados como classificadores. Serão avaliados diferentes números de qubits e estratégias de codificação dos dados, com o objetivo de analisar seus impactos no desempenho da classificação de imagens.
  \end{block}

  \vspace{0.25cm}

  \begin{center}
  \begin{tikzpicture}[
      box/.style={
        draw,
        rounded corners,
        align=center,
        minimum width=2.25cm,
        minimum height=0.8cm
      },
      arrow/.style={-{Latex[length=2mm]},thick}
    ]
    \node[box,fill=SoftGray] (img) {Imagem};
    \node[box,fill=NeuralOrange!20,right=0.35cm of img] (cnn) {Extrator de\\características};
    \node[box,fill=NeuralOrange!10,right=0.35cm of cnn] (feat) {Vetor de\\características};
    \node[box,fill=QuantumPurple!20,right=0.35cm of feat] (vqc) {Circuito quântico\\variacional};
    \node[box,fill=AccentGreen!20,right=0.35cm of vqc] (cls) {Classe\\predita};

    \draw[arrow] (img) -- (cnn);
    \draw[arrow] (cnn) -- (feat);
    \draw[arrow] (feat) -- (vqc);
    \draw[arrow] (vqc) -- (cls);
  \end{tikzpicture}
  \end{center}

  \begin{alertblock}{Pergunta orientadora}
    Como o número de qubits e a estratégia de codificação das características influenciam o desempenho de classificadores híbridos baseados em \highlight{circuitos quânticos variacionais} \cite{cerezo2021variational} na classificação de imagens, quando comparados a um baseline clássico treinado sobre as mesmas características?
  \end{alertblock}
\end{frame}

\begin{frame}{Motivação}
  \begin{columns}[T,onlytextwidth]
    
    \begin{column}{0.32\textwidth}
      \begin{block}{Aprendizado profundo}
        Redes neurais convolucionais pré-treinadas são capazes de extrair representações visuais discriminativas e reutilizáveis, reduzindo a necessidade de treinar todo o modelo a partir do zero.
      \end{block}
    \end{column}
    
    \begin{column}{0.32\textwidth}
      \begin{block}{Computação quântica}
        Circuitos quânticos parametrizados oferecem uma abordagem alternativa para codificar e transformar dados em espaços de Hilbert de alta dimensão \cite{schuld2015introduction,biamonte2017quantum}.
      \end{block}
    \end{column}
    
    \begin{column}{0.32\textwidth}
      \begin{block}{Arquitetura híbrida}
        A combinação de extratores clássicos com classificadores quânticos permite investigar o potencial dos circuitos variacionais sem transferir para o sistema quântico todo o custo do processamento visual.
      \end{block}
    \end{column}
    
  \end{columns}
\end{frame}

\begin{frame}{Pipeline de Referência: Classificador Clássico}
  \begin{center}
  \begin{tikzpicture}[
      box/.style={
        draw,
        rounded corners,
        align=center,
        minimum width=2.15cm,
        minimum height=0.9cm
      },
      arrow/.style={-{Latex[length=2mm]},thick}
    ]
    \node[box,fill=SoftGray] (img) {Imagem\\CIFAR-10};
    \node[box,fill=NeuralOrange!20,right=0.35cm of img] (resnet)
      {ResNet18\\pré-treinada};
    \node[box,fill=NeuralOrange!10,right=0.35cm of resnet] (feat)
      {Vetor de\\características};
    \node[box,fill=DeepBlue!10,right=0.35cm of feat] (proj)
      {Redução\\dimensional};
    \node[box,fill=DeepBlue!20,right=0.35cm of proj] (mlp)
      {MLP\\clássica};
    \node[box,fill=AccentGreen!20,right=0.35cm of mlp] (class)
      {Classe\\predita};

    \draw[arrow] (img) -- (resnet);
    \draw[arrow] (resnet) -- (feat);
    \draw[arrow] (feat) -- (proj);
    \draw[arrow] (proj) -- (mlp);
    \draw[arrow] (mlp) -- (class);
  \end{tikzpicture}
  \end{center}

  \vspace{0.3cm}

  \begin{block}{Baseline experimental}
    O desempenho da MLP será utilizado como referência para avaliar a
    competitividade dos classificadores quânticos. Para garantir uma comparação justa, os modelos clássico e quântico receberão as mesmas características extraídas pela ResNet18 escolhida como extratora de características.
  \end{block}
\end{frame}

\begin{frame}{Pipeline Proposto: Classificador Quântico}
  \begin{center}
  \begin{tikzpicture}[
      box/.style={
        draw,
        rounded corners,
        align=center,
        minimum width=2.05cm,
        minimum height=0.9cm
      },
      arrow/.style={-{Latex[length=2mm]},thick}
    ]
    \node[box,fill=SoftGray] (img) {Imagem\\CIFAR-10};
    \node[box,fill=NeuralOrange!20,right=0.3cm of img] (cnn)
      {ResNet18\\pré-treinada};
    \node[box,fill=NeuralOrange!10,right=0.3cm of cnn] (proj)
      {Redução\\dimensional};
    \node[box,fill=QuantumPurple!15,right=0.3cm of proj] (encode)
      {Codificação\\quântica};
    \node[box,fill=QuantumPurple!25,right=0.3cm of encode] (vqc)
      {Circuito quântico\\variacional};
    \node[box,fill=AccentGreen!20,right=0.3cm of vqc] (out)
      {Medição e\\predição};

    \draw[arrow] (img) -- (cnn);
    \draw[arrow] (cnn) -- (proj);
    \draw[arrow] (proj) -- (encode);
    \draw[arrow] (encode) -- (vqc);
    \draw[arrow] (vqc) -- (out);
  \end{tikzpicture}
  \end{center}

  \vspace{0.25cm}

  \begin{columns}[T,onlytextwidth]
    \begin{column}{0.48\textwidth}
      \begin{block}{Etapa clássica}
        A ResNet18 extrai representações visuais das imagens, enquanto a redução
        dimensional ajusta o vetor de características à quantidade de informações
        que pode ser codificada pelo circuito quântico.
      \end{block}
    \end{column}

    \begin{column}{0.48\textwidth}
      \begin{block}{Etapa quântica}
        O vetor reduzido é codificado em um estado quântico e processado por
        camadas parametrizadas. As medições realizadas ao final do circuito são
        utilizadas para produzir a predição da classe.
      \end{block}
    \end{column}
  \end{columns}

  \vspace{0.2cm}

  \begin{alertblock}{Ambiente de execução}
    Inicialmente, os circuitos quânticos serão executados em simuladores,
    permitindo avaliar diferentes configurações de qubits e estratégias de
    codificação em um ambiente controlado e livre de ruídos de hardware.
  \end{alertblock}
\end{frame}


\section{Fundamentação Metodológica}

%------------------------------------------------
\begin{frame}{Fundamentação: Qubits e Estados Quânticos}

  \begin{columns}[T,onlytextwidth]

    \begin{column}{0.48\textwidth}
      \begin{block}{Qubit}
        Diferentemente de um bit clássico, que assume os valores $0$ ou $1$,
        um qubit pode ser representado por uma superposição desses estados:
        \[
          \ket{\psi} = \alpha\ket{0} + \beta\ket{1},
        \]
        em que $\alpha,\beta\in\mathbb{C}$ e
        \[
          |\alpha|^2 + |\beta|^2 = 1.
        \]
      \end{block}
    \end{column}

    \begin{column}{0.48\textwidth}
      \begin{block}{Sistema com múltiplos qubits}
        Um sistema com $n$ qubits é descrito em um espaço de dimensão $2^n$:
        \[
          \ket{\psi}
          =
          \sum_{i=0}^{2^n-1} c_i\ket{i},
        \]
        com
        \[
          \sum_{i=0}^{2^n-1}|c_i|^2 = 1.
        \]
        Os coeficientes $c_i$ correspondem às amplitudes associadas aos
        estados da base computacional.
      \end{block}
    \end{column}

  \end{columns}
\end{frame}


%------------------------------------------------
\begin{frame}{Codificação de Dados em Estados Quânticos}

  \begin{alertblock}{Interpretação para aprendizado de máquina}
    Os dados clássicos devem ser transformados em operações ou amplitudes
    compatíveis com um estado quântico. Esse processo é denominado
    \highlight{codificação dos dados}.
  \end{alertblock}
  
  \begin{columns}[T,onlytextwidth]

    \begin{column}{0.48\textwidth}
      \begin{block}{Angle encoding}
        Cada característica controla o ângulo de uma porta de rotação:
        \[
          \ket{\psi(\mathbf{x})}
          =
          \bigotimes_{i=1}^{n}
          R_y(x_i)\ket{0}.
        \]

        \begin{itemize}
          \item Implementação simples;
          \item Baixa profundidade;
          \item Em geral, uma característica por qubit e por camada.
        \end{itemize}
      \end{block}
    \end{column}

    \begin{column}{0.48\textwidth}
      \begin{block}{Amplitude encoding}
        As características são inseridas nas amplitudes do estado:
        \[
          \ket{\psi(\mathbf{x})}
          =
          \sum_{i=0}^{2^n-1} x_i\ket{i}.
        \]

        \begin{itemize}
          \item Codifica até $2^n$ valores;
          \item Requer normalização;
          \item A preparação do estado pode ser custosa.
        \end{itemize}
      \end{block}
    \end{column}
  \end{columns}
\end{frame}


%------------------------------------------------
% \begin{frame}{Codificação Repetida: Data Re-uploading}

%   \begin{columns}[T,onlytextwidth]

%     \begin{column}{0.55\textwidth}
%       \begin{block}{Ideia}
%         Em vez de codificar os dados apenas no início do circuito,
%         o vetor de características é inserido novamente entre as camadas
%         variacionais:
%         \[
%           U(\mathbf{x},\boldsymbol{\theta})
%           =
%           \prod_{\ell=1}^{L}
%           U_{\mathrm{var}}(\boldsymbol{\theta}_{\ell})
%           U_{\mathrm{enc}}(\mathbf{x}).
%         \]

%         Essa estratégia pode aumentar a expressividade do modelo mesmo
%         quando poucos qubits estão disponíveis.
%       \end{block}
%     \end{column}

%     \begin{column}{0.41\textwidth}
%       \begin{center}
%       \begin{tikzpicture}[
%           box/.style={
%             draw,
%             rounded corners,
%             align=center,
%             minimum width=2.6cm,
%             minimum height=0.7cm
%           },
%           arrow/.style={-{Latex[length=2mm]},thick},
%           node distance=0.3cm
%         ]
%         \node[box,fill=QuantumPurple!15] (enc1)
%           {Codificação\\$U_{\mathrm{enc}}(\mathbf{x})$};

%         \node[box,fill=QuantumPurple!25,below=of enc1] (var1)
%           {Camada variacional\\$U_{\mathrm{var}}(\theta_1)$};

%         \node[box,fill=QuantumPurple!15,below=of var1] (enc2)
%           {Nova codificação\\$U_{\mathrm{enc}}(\mathbf{x})$};

%         \node[box,fill=QuantumPurple!25,below=of enc2] (var2)
%           {Camada variacional\\$U_{\mathrm{var}}(\theta_2)$};

%         \node[box,fill=AccentGreen!20,below=of var2] (med)
%           {Medição};

%         \draw[arrow] (enc1) -- (var1);
%         \draw[arrow] (var1) -- (enc2);
%         \draw[arrow] (enc2) -- (var2);
%         \draw[arrow] (var2) -- (med);
%       \end{tikzpicture}
%       \end{center}
%     \end{column}

%   \end{columns}

%   \begin{alertblock}{Hipótese experimental}
%     O data re-uploading pode compensar parcialmente a limitação no número
%     de qubits, mas também aumenta a profundidade e o custo do circuito.
%   \end{alertblock}

% \end{frame}


%------------------------------------------------
\begin{frame}{Circuito Quântico Variacional}

  \begin{columns}[T,onlytextwidth]

    \begin{column}{0.48\textwidth}
      \begin{block}{Estrutura}
        Um circuito quântico variacional combina:
        \begin{itemize}
          \item uma etapa de codificação dos dados;
          \item portas parametrizadas;
          \item operações de emaranhamento;
          \item medições ao final do circuito.
        \end{itemize}

        Sua saída pode ser expressa como:
        \[
          f_{\boldsymbol{\theta}}(\mathbf{x})
          =
          \bra{\psi(\mathbf{x})}
          U^\dagger(\boldsymbol{\theta})
          O
          U(\boldsymbol{\theta})
          \ket{\psi(\mathbf{x})},
        \]
        em que $O$ é um observável quântico.
      \end{block}
    \end{column}

    \begin{column}{0.48\textwidth}
      \begin{block}{Ansatz variacional}
        As camadas podem ser compostas por rotações parametrizadas e portas
        de emaranhamento:
        \[
          U_{\mathrm{var}}(\boldsymbol{\theta})
          =
          U_{\mathrm{ent}}
          \prod_{i=1}^{n}
          R_z(\theta_{i,3})
          R_y(\theta_{i,2})
          R_x(\theta_{i,1}).
        \]

        \begin{itemize}
          \item As rotações possuem parâmetros treináveis;
          \item As portas CNOT promovem emaranhamento;
          \item O bloco pode ser repetido por $L$ camadas.
        \end{itemize}
      \end{block}
    \end{column}

  \end{columns}
\end{frame}

%------------------------------------------------
\begin{frame}{Treinamento do Modelo Híbrido}

  \begin{center}
  \begin{tikzpicture}[
      box/.style={
        draw,
        rounded corners,
        align=center,
        minimum width=2.1cm,
        minimum height=0.8cm
      },
      arrow/.style={-{Latex[length=2mm]},thick}
    ]
    \node[box,fill=SoftGray] (input) {Características\\reduzidas};
    \node[box,fill=QuantumPurple!15,right=0.35cm of input] (enc)
      {Codificação\\quântica};
    \node[box,fill=QuantumPurple!25,right=0.35cm of enc] (vqc)
      {Circuito\\variacional};
    \node[box,fill=AccentGreen!15,right=0.35cm of vqc] (med)
      {Valores\\esperados};
    \node[box,fill=DeepBlue!15,right=0.35cm of med] (loss)
      {Função de\\perda};

    \draw[arrow] (input) -- (enc);
    \draw[arrow] (enc) -- (vqc);
    \draw[arrow] (vqc) -- (med);
    \draw[arrow] (med) -- (loss);

    \draw[arrow,bend left=35]
      (loss.north) to node[above,align=center]
      {\scriptsize atualização dos\\[-0.1cm]\scriptsize parâmetros}
      (vqc.north);
  \end{tikzpicture}
  \end{center}

  \vspace{0.3cm}

  \begin{columns}[T,onlytextwidth]

    \begin{column}{0.48\textwidth}
      \begin{block}{Predição}
        Para cada qubit medido, pode-se calcular o valor esperado do
        observável de Pauli-$Z$:
        \[
          z_i = \langle Z_i\rangle.
        \]

        O vetor
        \[
          \mathbf{z} =
          [\langle Z_1\rangle,\ldots,\langle Z_m\rangle]
        \]
        é utilizado para produzir a predição da classe.
      \end{block}
    \end{column}

    \begin{column}{0.48\textwidth}
      \begin{block}{Otimização}
        A função de perda é calculada de forma clássica, e seus gradientes
        são utilizados para atualizar os parâmetros do circuito:
        \[
          \boldsymbol{\theta}
          \leftarrow
          \boldsymbol{\theta}
          -
          \eta\nabla_{\boldsymbol{\theta}}\mathcal{L}.
        \]

        O gradiente pode ser obtido por métodos como a regra
        \textit{parameter-shift}.
      \end{block}
    \end{column}

  \end{columns}

\end{frame}

\begin{frame}{Papel do Numero de Qubits}
  \begin{columns}[T,onlytextwidth]
    \begin{column}{0.48\textwidth}
      \begin{block}{Hipotese}
        Mais qubits aumentam a capacidade de representacao do circuito, mas tambem aumentam custo de simulacao, risco de overfitting e dificuldade de otimizacao.
      \end{block}
      \begin{itemize}
        \item Poucos qubits: classificador mais restrito.
        \item Mais qubits: maior espaco de estados.
        \item Simulacao classica cresce rapidamente.
      \end{itemize}
    \end{column}
    \begin{column}{0.48\textwidth}
      \begin{center}
      \begin{tikzpicture}
        \draw[->,thick] (0,0) -- (4.4,0) node[right] {qubits};
        \draw[->,thick] (0,0) -- (0,3.0) node[above] {custo};
        \draw[QuantumPurple,very thick] plot[smooth] coordinates {(0.3,0.3) (1.1,0.55) (2.0,1.0) (2.8,1.75) (3.7,2.8)};
        \node[below] at (0.8,0) {4};
        \node[below] at (1.8,0) {6};
        \node[below] at (2.8,0) {8};
        \node[below] at (3.8,0) {10};
        \node[align=center,text width=3.4cm] at (2.4,2.5) {Trade-off entre capacidade e viabilidade};
      \end{tikzpicture}
      \end{center}
    \end{column}
  \end{columns}
\end{frame}

\section{Metodologia}

\begin{frame}{O que será comparado?}
  \begin{columns}[T,onlytextwidth]

    \begin{column}{0.32\textwidth}
      \begin{block}{Codificação dos dados}
        \begin{itemize}
          \item Angle encoding
          \item Amplitude encoding
          \item Basis encoding
          \item Data re-uploading
        \end{itemize}
      \end{block}
    \end{column}

    \begin{column}{0.32\textwidth}
      \begin{block}{Número de qubits}
        \begin{itemize}
          \item 4 qubits
          \item 6 qubits
          \item 8 qubits
          \item 10 qubits
        \end{itemize}
      \end{block}
    \end{column}

    \begin{column}{0.32\textwidth}
      \begin{block}{Circuito variacional}
        \begin{itemize}
          \item Número de camadas
          \item Padrão de emaranhamento
          \item Número de parâmetros
          \item Profundidade do circuito
        \end{itemize}
      \end{block}
    \end{column}

  \end{columns}

  \vspace{0.2cm}

  \begin{alertblock}{Estratégia experimental}
    Cada fator será analisado de forma controlada, mantendo os demais
    componentes fixos, de modo a identificar sua influência individual
    sobre o desempenho do classificador.
  \end{alertblock}
\end{frame}

\begin{frame}{Desenho Experimental}
  \begin{enumerate}
    \item Utilizar uma ResNet18 pré-treinada no ImageNet como extratora de características.
    
    \item Manter os pesos da ResNet18 congelados durante os experimentos, garantindo que todos os classificadores utilizem o mesmo espaço de características.
    
    \item Aplicar um método de redução ou projeção dimensional para adequar o vetor de características ao número de qubits e à estratégia de codificação avaliados.
    
    \item Treinar circuitos quânticos variacionais para diferentes combinações de número de qubits, forma de codificação e arquitetura do circuito.
    
    \item Treinar uma MLP clássica como baseline, utilizando as mesmas características reduzidas, os mesmos conjuntos de treino, validação e teste e condições experimentais equivalentes.
    
    \item Comparar os modelos com base no desempenho preditivo, no custo computacional, no número de parâmetros e na estabilidade entre diferentes execuções.
  \end{enumerate}

  \vspace{0.25cm}

  \begin{alertblock}{Controle experimental}
    Para garantir uma comparação justa, todos os modelos utilizarão os mesmos dados, a mesma divisão entre treino, validação e teste e as mesmas características extraídas pela rede convolucional antes da etapa de classificação.
  \end{alertblock}
\end{frame}

\begin{frame}{Metricas e Criterios de Analise}
  \begin{columns}[T,onlytextwidth]
    \begin{column}{0.48\textwidth}
      \begin{block}{Desempenho preditivo}
        \begin{itemize}
          \item Acuracia
          \item Loss de validacao/teste
          \item Matriz de confusao
          \item Media e desvio entre folds
        \end{itemize}
      \end{block}
    \end{column}
    \begin{column}{0.48\textwidth}
      \begin{block}{Custo e estabilidade}
        \begin{itemize}
          \item Tempo de treinamento
          \item Numero de parametros
          \item Numero de qubits
          \item Profundidade do circuito
          \item Variabilidade entre execucoes
        \end{itemize}
      \end{block}
    \end{column}
  \end{columns}
\end{frame}

\begin{frame}{Hipoteses de Pesquisa}
  \begin{block}{H1: codificacao importa}
    Diferentes estrategias de codificacao podem produzir desempenhos distintos mesmo quando o extrator classico e mantido fixo.
  \end{block}
  \begin{block}{H2: mais qubits nem sempre e melhor}
    O aumento do numero de qubits pode melhorar expressividade, mas tambem tornar a otimizacao mais dificil e aumentar custo computacional.
  \end{block}
  \begin{block}{H3: modelo hibrido pode ser competitivo}
    Um classificador quantico variacional pode se aproximar da MLP classica em alguns cenarios, especialmente com caracteristicas ja bem estruturadas.
  \end{block}
\end{frame}

\begin{frame}{Etapas de Implementacao}
  \begin{center}
  \begin{tikzpicture}[
      flowstep/.style={draw,rounded corners,fill=SoftGray,align=center,minimum width=2.3cm,minimum height=0.85cm},
      arrow/.style={-{Latex[length=2mm]},thick}
    ]
    \node[flowstep] (s1) {Baseline\\classico};
    \node[flowstep,right=0.35cm of s1] (s2) {Extracao de\\features};
    \node[flowstep,right=0.35cm of s2] (s3) {Modulo\\quantico};
    \node[flowstep,right=0.35cm of s3] (s4) {Grade de\\experimentos};
    \node[flowstep,right=0.35cm of s4] (s5) {Analise e\\relatorio};
    \draw[arrow] (s1) -- (s2);
    \draw[arrow] (s2) -- (s3);
    \draw[arrow] (s3) -- (s4);
    \draw[arrow] (s4) -- (s5);
  \end{tikzpicture}
  \end{center}

  \vspace{0.4cm}
  \begin{itemize}
    \item Integrar biblioteca quantica: PennyLane \cite{bergholm2018pennylane} ou Qiskit Machine Learning.
    \item Criar interface comum para classificadores classico e quantico.
    \item Automatizar experimentos por codificacao, qubits e profundidade.
  \end{itemize}
\end{frame}

\begin{frame}{Riscos e Limitacoes}
  \begin{columns}[T,onlytextwidth]
    \begin{column}{0.48\textwidth}
      \begin{alertblock}{Computacionais}
        \begin{itemize}
          \item Simulacao quantica pode ser lenta.
          \item Muitos qubits podem inviabilizar a grade completa.
          \item Amplitude encoding pode ter alto custo de preparacao.
        \end{itemize}
      \end{alertblock}
    \end{column}
    \begin{column}{0.48\textwidth}
      \begin{alertblock}{Metodologicos}
        \begin{itemize}
          \item Comparacao precisa controlar splits e features.
          \item Resultados podem depender do ansatz escolhido.
          \item Ganho sobre MLP classica nao e garantido.
        \end{itemize}
      \end{alertblock}
    \end{column}
  \end{columns}
\end{frame}

\begin{frame}{Resultado Esperado da Proposta}
  \begin{block}{Produto tecnico}
    Um pipeline experimental hibrido capaz de comparar classificadores classicos e quanticos sobre as mesmas caracteristicas profundas extraidas de imagens.
  \end{block}

  \begin{block}{Produto cientifico}
    Uma analise empirica sobre como a escolha da codificacao quantica e do numero de qubits afeta desempenho, custo e estabilidade em classificacao de imagens.
  \end{block}

  \begin{block}{Entrega final}
    Codigo reprodutivel, tabelas comparativas, graficos de desempenho e discussao das limitacoes observadas.
  \end{block}
\end{frame}

\begin{frame}{Referencias Iniciais}
  \small
  \bibliographystyle{plainnat}
  \bibliography{referencias}
\end{frame}

\end{document}
